% ============================================================================
% Chapter 8 solutions: Simplex - Matrix Calculations
% Book source: part1-linear-programming/ch06-simplex/simplex-matrix-version.tex
% Exercise numbering follows \begin{ex} order in that file (9 exercises, all
% in the Exercises section; no embedded mid-chapter exercise exists in the
% current source).
% All basis computations verified with exact-fraction linear algebra and
% scipy.optimize.linprog, July 2026.
% ============================================================================
\section{Chapter 8: Simplex - Matrix Calculations}

Chapter 8 recasts the simplex method as basis algebra: $A_B$, $A_B^{-1}\b$,
and the reduced costs $\c_N^\top - \c_B^\top A_B^{-1} A_N$. Its nine
exercises split into three warm-ups on computing basic solutions from a
basis (8.1 to 8.3), three core computations with reduced costs and the
dictionary formula (8.4 to 8.6), two conceptual exercises connecting
reduced costs to the dictionary and tableau sign conventions (8.7, 8.8),
and one challenge proof on uniqueness of the optimum (8.9). Every matrix
inverse, basic solution, and reduced-cost vector below was verified with
exact-fraction computations. Book Selected Solutions exist for 8.1, 8.4,
and 8.9 and are reproduced here.

\exsol{8.1}{Basis matrix and basic solution}{1}

(Book Selected Solution, verified.) The columns of $x_2$ and $x_3$ give
\[
A_B =
\begin{bmatrix}
2 & 1 \\
1 & 2
\end{bmatrix},
\qquad
A_B^{-1} = \frac{1}{3}
\begin{bmatrix}
2 & -1 \\
-1 & 2
\end{bmatrix}.
\]
The basic solution is
\[
\x_B = A_B^{-1}\b = \frac{1}{3}
\begin{bmatrix}
2 & -1 \\
-1 & 2
\end{bmatrix}
\begin{bmatrix}
4 \\ 5
\end{bmatrix}
= \frac{1}{3}
\begin{bmatrix}
3 \\ 6
\end{bmatrix}
=
\begin{bmatrix}
1 \\ 2
\end{bmatrix},
\]
so $(x_1, x_2, x_3) = (0, 1, 2)$. Both basic variables are nonnegative,
so this is a basic feasible solution.

\exsol{8.2}{A basic solution that is not feasible}{1}

The columns of $x_1$ and $x_3$ give
\[
A_B =
\begin{bmatrix}
1 & 1 \\
3 & 2
\end{bmatrix},
\qquad
\det A_B = 2 - 3 = -1,
\qquad
A_B^{-1} =
\begin{bmatrix}
-2 & 1 \\
3 & -1
\end{bmatrix}.
\]
Setting $x_2 = 0$, the basic solution is
\[
\x_B = A_B^{-1}\b =
\begin{bmatrix}
-2 & 1 \\
3 & -1
\end{bmatrix}
\begin{bmatrix}
4 \\ 5
\end{bmatrix}
=
\begin{bmatrix}
-3 \\ 7
\end{bmatrix},
\]
so $(x_1, x_2, x_3) = (-3, 0, 7)$. It solves $A\x = \b$ (check:
$-3 + 2(0) + 7 = 4$ and $3(-3) + 0 + 2(7) = 5$), but the basic variable
$x_1 = -3$ violates the nonnegativity constraint $\x \geq \mathbf{0}$. A basic
solution only qualifies as a basic \emph{feasible} solution when
$\x_B \geq \mathbf{0}$, so this basis gives an infeasible basic solution.

\exsol{8.3}{A fourth basis for the three-bases example}{1}

The columns of $x_2$ and $s_2$ give
\[
A_B =
\begin{bmatrix}
2 & 0 \\
3 & 1
\end{bmatrix},
\qquad
A_B^{-1} = \frac{1}{2}
\begin{bmatrix}
1 & 0 \\
-3 & 2
\end{bmatrix}
=
\begin{bmatrix}
\tfrac12 & 0 \\
-\tfrac32 & 1
\end{bmatrix}.
\]
The basic solution is
\[
\x_B =
\begin{bmatrix} x_2 \\ s_2 \end{bmatrix}
= A_B^{-1}\b =
\begin{bmatrix}
\tfrac12 & 0 \\
-\tfrac32 & 1
\end{bmatrix}
\begin{bmatrix}
16 \\ 45
\end{bmatrix}
=
\begin{bmatrix}
8 \\ 21
\end{bmatrix}.
\]
Both components are nonnegative, so the basis is \emph{feasible}. Since
$x_1 = s_1 = 0$, the corresponding point is $(x_1, x_2) = (0, 8)$: the
vertex where the boundary $x_1 + 2x_2 = 16$ meets the $x_2$-axis. (This
is the basis shown to be optimal in the worked example of the optimality
section.)

\exsol{8.4}{Reduced costs at two bases}{2}

(Book Selected Solution, verified, with the intermediate matrices shown.)

\textbf{Basis $B_1 = \{y, s_1, s_2\}$.} Ordering the basic variables
$(y, s_1, s_2)$,
\[
A_{B_1} =
\begin{bmatrix}
1 & 1 & 0 \\
1 & 0 & 1 \\
2 & 0 & 0
\end{bmatrix},
\qquad
A_{B_1}^{-1} =
\begin{bmatrix}
0 & 0 & \tfrac12 \\
1 & 0 & -\tfrac12 \\
0 & 1 & -\tfrac12
\end{bmatrix},
\qquad
\x_{B_1} = A_{B_1}^{-1}\b = \begin{bmatrix} 7 \\ 2 \\ 9 \end{bmatrix},
\]
so $y = 7$, $s_1 = 2$, $s_2 = 9$: a basic feasible solution with objective
value $\c_{B_1}^\top \x_{B_1} = 3 \cdot 7 = 21$. With
$\c_{B_1}^\top = (3, 0, 0)$, the simplex multipliers are
$\c_{B_1}^\top A_{B_1}^{-1} = (0, 0, \tfrac32)$, and for the nonbasic
variables $x$ and $s_3$ (columns $(1,2,1)^\top$ and $(0,0,1)^\top$),
\[
\c_N^\top - \c_{B_1}^\top A_{B_1}^{-1} A_N
= \begin{bmatrix} 2 & 0 \end{bmatrix}
- \begin{bmatrix} \tfrac32 & \tfrac32 \end{bmatrix}
= \begin{bmatrix} \tfrac12 & -\tfrac32 \end{bmatrix}.
\]
The reduced cost of $x$ is $\tfrac12 > 0$, so $B_1$ is \emph{not}
optimal; this matches the second dictionary of Chapter~7,
$z = 21 + \tfrac12 x - \tfrac32 s_3$.

\textbf{Basis $B_2 = \{x, y, s_2\}$.} Ordering the basic variables
$(x, y, s_2)$,
\[
A_{B_2} =
\begin{bmatrix}
1 & 1 & 0 \\
2 & 1 & 1 \\
1 & 2 & 0
\end{bmatrix},
\qquad
A_{B_2}^{-1} =
\begin{bmatrix}
2 & 0 & -1 \\
-1 & 0 & 1 \\
-3 & 1 & 1
\end{bmatrix},
\qquad
\x_{B_2} = A_{B_2}^{-1}\b = \begin{bmatrix} 4 \\ 5 \\ 3 \end{bmatrix},
\]
a basic feasible solution with value $2 \cdot 4 + 3 \cdot 5 = 23$. With
$\c_{B_2}^\top = (2, 3, 0)$ the multipliers are
$\c_{B_2}^\top A_{B_2}^{-1} = (1, 0, 1)$, and for the nonbasic variables
$s_1, s_3$ (columns $(1,0,0)^\top$ and $(0,0,1)^\top$),
\[
\c_N^\top - \c_{B_2}^\top A_{B_2}^{-1} A_N
= \begin{bmatrix} 0 & 0 \end{bmatrix}
- \begin{bmatrix} 1 & 1 \end{bmatrix}
= \begin{bmatrix} -1 & -1 \end{bmatrix} \leq \mathbf{0}^\top,
\]
so $B_2$ is optimal, matching the final dictionary
$z = 23 - s_1 - s_3$.

\exsol{8.5}{Which bases are feasible?}{2}

\textbf{(a)} The six column pairs of
$A = \begin{bmatrix} 1 & 2 & 1 & 0 \\ 5 & 3 & 0 & 1 \end{bmatrix}$ have
determinants
\[
\begin{aligned}
&\det A_{\{x_1,x_2\}} = -7,\qquad
\det A_{\{x_1,s_1\}} = -5,\qquad
\det A_{\{x_1,s_2\}} = 1,\\
&\det A_{\{x_2,s_1\}} = -3,\qquad
\det A_{\{x_2,s_2\}} = 2,\qquad
\det A_{\{s_1,s_2\}} = 1,
\end{aligned}
\]
all nonzero, so every pair of columns is linearly independent and each
choice is a basis.

\textbf{(b)} Solving $A_B \x_B = \b$ for each basis (nonbasic variables
zero):
\begin{center}
\begin{tabular}{l|l|l|c}
Basis & Basic solution & Point $(x_1, x_2)$ & Feasible? \\
\hline
$\{x_1, x_2\}$ & $x_1 = 6,\ x_2 = 5$ & $(6, 5)$ & yes \\
$\{x_1, s_1\}$ & $x_1 = 9,\ s_1 = 7$ & $(9, 0)$ & yes \\
$\{x_1, s_2\}$ & $x_1 = 16,\ s_2 = -35$ & $(16, 0)$ & no ($s_2 < 0$) \\
$\{x_2, s_1\}$ & $x_2 = 15,\ s_1 = -14$ & $(0, 15)$ & no ($s_1 < 0$) \\
$\{x_2, s_2\}$ & $x_2 = 8,\ s_2 = 21$ & $(0, 8)$ & yes \\
$\{s_1, s_2\}$ & $s_1 = 16,\ s_2 = 45$ & $(0, 0)$ & yes \\
\end{tabular}
\end{center}
The four feasible bases correspond to the four vertices of the feasible
region, $(0,0)$, $(9,0)$, $(6,5)$, $(0,8)$; the two infeasible bases are
intersections of constraint boundaries lying outside the region.

\textbf{(c)} Take
\[
A = \begin{bmatrix} 1 & 2 & 1 & 0 \\ 2 & 4 & 0 & 1 \end{bmatrix}.
\]
The first two columns satisfy $(2, 4)^\top = 2 \cdot (1, 2)^\top$, so
the pair $\{1, 2\}$ is linearly dependent ($\det = 0$) and does not form
a basis.

\exsol{8.6}{Rebuilding a dictionary from the formula}{2}

With variables ordered $(x, y, s_1, s_2)$,
\[
A = \begin{bmatrix} 3 & 1 & 1 & 0 \\ 5 & 3 & 0 & 1 \end{bmatrix},
\quad
\b = \begin{bmatrix} 23 \\ 45 \end{bmatrix},
\quad
\c^\top = \begin{bmatrix} 2 & 1 & 0 & 0 \end{bmatrix},
\quad
B = \{x, y\},\ N = \{s_1, s_2\}.
\]
Then
\[
A_B = \begin{bmatrix} 3 & 1 \\ 5 & 3 \end{bmatrix},
\qquad
A_B^{-1} = \frac{1}{4} \begin{bmatrix} 3 & -1 \\ -5 & 3 \end{bmatrix}
= \begin{bmatrix} 0.75 & -0.25 \\ -1.25 & 0.75 \end{bmatrix}.
\]
Since $A_N = I_2$ (the slack columns),
\[
A_B^{-1}\b = \frac14 \begin{bmatrix} 3(23) - 45 \\ -5(23) + 3(45) \end{bmatrix}
= \begin{bmatrix} 6 \\ 5 \end{bmatrix},
\qquad
A_B^{-1}A_N = A_B^{-1}
= \begin{bmatrix} 0.75 & -0.25 \\ -1.25 & 0.75 \end{bmatrix}.
\]
The constraint part of the dictionary,
$\x_B = A_B^{-1}\b - A_B^{-1}A_N \x_N$, reads
\[
\begin{aligned}
x &= 6 - 0.75\,s_1 + 0.25\,s_2, \\
y &= 5 + 1.25\,s_1 - 0.75\,s_2,
\end{aligned}
\]
exactly the displayed rows. For the objective,
$\c_B^\top A_B^{-1}\b = 2(6) + 1(5) = 17$, and
\[
\c_B^\top A_B^{-1} = \begin{bmatrix} 2 & 1 \end{bmatrix}
\begin{bmatrix} 0.75 & -0.25 \\ -1.25 & 0.75 \end{bmatrix}
= \begin{bmatrix} 0.25 & 0.25 \end{bmatrix},
\qquad
\c_N^\top - \c_B^\top A_B^{-1}A_N
= \begin{bmatrix} -0.25 & -0.25 \end{bmatrix}.
\]
So the formula gives
$\max\ 17 - 0.25 s_1 - 0.25 s_2$, reproducing the displayed dictionary
in full. (Both reduced costs are negative, so this basis is optimal.)

\exsol{8.7}{Reduced costs of basic variables}{2}

Applying the reduced-cost formula to the basic columns replaces $A_N$ by
$A_B$:
\[
\c_B^\top - \c_B^\top A_B^{-1} A_B
= \c_B^\top - \c_B^\top I_m
= \c_B^\top - \c_B^\top = \mathbf{0}^\top,
\]
since $A_B^{-1}A_B = I_m$ by the definition of the inverse. So every
basic variable has reduced cost exactly zero.

Interpretation: the dictionary at basis $B$ expresses $z$ and each basic
variable purely in terms of the \emph{nonbasic} variables. Constructing
the dictionary means substituting
$\x_B = A_B^{-1}\b - A_B^{-1}A_N\x_N$ into the objective, which
eliminates every basic variable from the objective row; the computation
above is the algebraic statement that this elimination is exact. A basic
variable can never appear in the objective row because, if it did, the
row would not be solved for the basic variables, and the ``dictionary''
property (basic variables and $z$ as functions of nonbasic variables
alone) would fail. Equivalently: the value of a basic variable is
determined by the nonbasic variables, so it carries no independent rate
of change for $z$.

\exsol{8.8}{Reduced-cost signs and the entering rule}{2}

\textbf{(a)} Rewriting the objective from the viewpoint of basis $B$,
\[
z = \c_B^\top A_B^{-1}\b
+ \left(\c_N^\top - \c_B^\top A_B^{-1} A_N\right)\x_N,
\]
which is precisely the objective row of the dictionary: a constant plus a
linear function of the nonbasic variables. The coefficient of a single
nonbasic variable $x_j$ in this row is the $j$th component of the
reduced-cost vector, namely $c_j - \c_B^\top A_B^{-1} A_{\{j\}}$. So
``pick a nonbasic variable with a positive objective-row coefficient''
and ``pick a nonbasic variable with a positive reduced cost'' are
literally the same rule, and the stopping conditions (no positive
coefficient; all reduced costs $\leq 0$) coincide as well.

\textbf{(b)} Increasing a nonbasic variable $x_j$ from zero, while
keeping the other nonbasic variables at zero, moves the solution along
the ray
\[
\x(t) = \bar\x + t\,\mathbf{d},
\qquad
\mathbf{d}_B = -A_B^{-1}A_{\{j\}},\quad d_j = 1,\quad d_k = 0
\text{ otherwise}.
\]
Geometrically this is a move away from the current vertex along an edge
of the feasible region: one tight constraint ($x_j = 0$) is released
while the equality constraints stay satisfied. The objective changes at
rate equal to the reduced cost, so a positive reduced cost means the edge
points uphill. The ratio test stops the move at the neighboring vertex,
where some basic variable reaches zero.

\textbf{(c)} The tableau stores the objective as the equation
$z - \c^\top\x = 0$, so the entry recorded under a nonbasic variable is
the \emph{negative} of its dictionary coefficient, i.e., the negative of
its reduced cost. Hunting for a negative $z$-row entry in the tableau is
therefore hunting for a positive reduced cost, and ``stop when all
$z$-row entries are nonnegative'' is ``stop when all reduced costs are
$\leq 0$.'' The three formulations state one test with two sign
conventions; no contradiction arises.

\exsol{8.9}{Strictly negative reduced costs imply a unique optimum}{3}

(Book Selected Solution, verified.) Write
$\bar{\c}_N^\top = \c_N^\top - \c_B^\top A_B^{-1} A_N$ and
$z^* = \c^\top \bar\x = \c_B^\top A_B^{-1} \b$. Every feasible $\x$
satisfies $A_B \x_B + A_N \x_N = \b$, so
$\x_B = A_B^{-1}\b - A_B^{-1}A_N \x_N$ and, exactly as in the optimality
theorem,
\[
\c^\top \x = z^* + \bar{\c}_N^\top \x_N .
\]
Take any feasible $\x \neq \bar\x$.

\emph{Case 1: $\x_N \neq \mathbf{0}$.} Then some component $x_j > 0$ with
$j \in N$, and since every entry of $\bar{\c}_N$ is strictly negative,
$\bar{\c}_N^\top \x_N < 0$. Hence $\c^\top \x < z^*$, so $\x$ is not
optimal.

\emph{Case 2: $\x_N = \mathbf{0}$.} Then $A_B \x_B = \b$, and since $A_B$ is
invertible this forces $\x_B = A_B^{-1}\b$, i.e., $\x = \bar\x$,
contradicting $\x \neq \bar\x$.

Therefore every feasible point other than $\bar\x$ has objective value
strictly less than $z^*$: the solution $\bar\x$ is optimal and no other
optimal solution exists.
